\chapter{系统软件设计}

\section{整体工作流程}

本节填写系统软件的整体工作流程。可说明程序从上电初始化到进入主循环后的主要执行过程，包括系统初始化、输入采集、数据处理、状态判断、输出控制和显示刷新等步骤。

实际撰写时，可结合具体课题绘制整体流程图，并在正文中说明各流程节点的作用。

\section{各模块设计}

本节填写系统软件中各个功能模块的设计内容。可根据课程设计实际功能，将软件划分为若干模块分别说明。

\subsection{输入检测模块}

本节填写输入检测模块的设计内容。可说明输入信号的来源、采集方式、判断条件以及与后续处理模块之间的关系。

\subsection{数据处理模块}

本节填写数据处理模块的设计内容。可说明系统如何对采集到的数据进行滤波、转换、计算、判断或状态更新。


\section{本章小结}

本章主要介绍了系统软件设计内容，包括整体工作流程、各功能模块设计等。通过本章设计，可以明确系统软件的实现思路，为后续系统测试与结果分析提供基础。